% ================= PART II =================
\part{The Mathematics Behind GLMs}

\chapter{The Exponential Family}

\Exploration
\tbc

\Cognition

\section{Definition and Properties}

A probability distribution belongs to the exponential family if its density (or mass) function can be written in the general form
\begin{mathbox}
\[
f_X(x\mid\theta) = h(x)\exp\left\{\eta(\theta)\cdot T(x) - A(\theta)\right\}
\]
\end{mathbox}
Here $h(x)$ is a known, non-negative function of the data alone, while $\eta(\theta)$, $T(x)$ and $A(\theta)$ are known functions governing the natural parameter, the sufficient statistic, and the normalising term respectively \citep{mccullagh1989}. In this general version $\theta$ may be a vector and $T(x)$ a vector of sufficient statistics, so this is sometimes called the multi-parameter or full exponential family.

For GLMs, though, what matters is the one-parameter exponential family, since this is the restricted form underlying the standard GLM distributions: Binomial, Poisson, Normal, and Gamma. Restricting $\theta$ to a scalar and writing the exponent as a sum rather than a dot product gives the working form used throughout this portfolio:
\begin{mathbox}
\[
f(y;\theta) = \exp\left\{a(y)\,b(\theta) + c(\theta) + d(y)\right\}
\]
\end{mathbox}
Matching the terms defined in Part I:
\begin{itemize}
\item $a(y)$ is the sufficient statistic: the function of the data that gets multiplied by the natural parameter.
\item $b(\theta)$ is the natural (canonical) parameter, the function of $\theta$ through which the linear predictor will eventually connect to the distribution.
\item $c(\theta)$ is the log-partition function, which guarantees $f(y;\theta)$ integrates (or sums) to one over its support.
\item $d(y)$ is the base measure, the leftover piece depending on $y$ alone, with no $\theta$ in it.
\end{itemize}
What makes this form useful is the separation it forces between data and parameter. $y$ and $\theta$ never appear multiplied together except through the product $a(y)\,b(\theta)$; everything else involving $y$ sits in $d(y)$, and everything else involving $\theta$ sits in $c(\theta)$. This separation is also what makes $a(y)$ sufficient in the first place: once $a(y)$ is known, nothing more about $y$ is needed to work out how the likelihood depends on $\theta$.

A distribution is in canonical form when $a(y)=y$. In that case $b(\theta)$ is called the canonical parameter, and, as covered later when link functions are discussed, the link function that sets the linear predictor equal to $b(\theta)$ directly is the canonical link for that distribution.

$c(\theta)$ is worth dwelling on a little longer, because it isn't really a free choice, it's forced by the requirement that $f(y;\theta)$ be a valid density. Given $a(y)$, $b(\theta)$ and $d(y)$, normalisation requires
\begin{mathbox}
\[
\int \exp\{a(y)\,b(\theta)+d(y)\}\,dy = \exp\{-c(\theta)\}
\]
\end{mathbox}
(or the equivalent sum in the discrete case), so $c(\theta)$ ends up completely determined by the rest of the exponent. This is also why $c(\theta)$ is called the cumulant function: differentiating it with respect to $\theta$ generates the moments of $a(y)$. Under the canonical parameterisation, $\mathbb{E}[a(Y)] = -c'(\theta)$ and $\operatorname{Var}[a(Y)] = -c''(\theta)$, a fact that becomes useful once the mean-variance relationships for individual GLM distributions are worked out further on.

\subsection*{Proof: $c(\theta)$ as the Cumulant (Normalising) Function}
The claim above, that $c(\theta)$ is forced by normalisation rather than being a free choice, is not asserted without justification; it follows directly from requiring $f(y;\theta)$ to be a valid density. The derivation below is a proof of that result.

Since $f(y;\theta)$ is a probability density (or mass) function, it must integrate, or sum in the discrete case, to 1:
\begin{mathbox}
\[
\int f(y;\theta)\,dy = 1
\]
\end{mathbox}
Substituting the one-parameter exponential family form gives
\begin{mathbox}
\[
\int \exp\{a(y)b(\theta)+c(\theta)+d(y)\}\,dy = 1
\]
\end{mathbox}
Since $c(\theta)$ does not depend on $y$, it can be factored out of the integral using $\exp\{A+B\}=\exp\{A\}\cdot\exp\{B\}$:
\begin{mathbox}
\[
\exp\{c(\theta)\}\int \exp\{a(y)b(\theta)+d(y)\}\,dy = 1
\]
\end{mathbox}
Dividing both sides by the remaining integral isolates $\exp\{c(\theta)\}$:
\begin{mathbox}
\[
\exp\{c(\theta)\} = \left(\int \exp\{a(y)b(\theta)+d(y)\}\,dy\right)^{-1}
\]
\end{mathbox}
Taking the logarithm of both sides,
\begin{mathbox}
\[
\log\left[\exp\{c(\theta)\}\right] = \log\left[\left(\int \exp\{a(y)b(\theta)+d(y)\}\,dy\right)^{-1}\right]
\]
\end{mathbox}
which simplifies, using $\log(x^{-1}) = -\log(x)$, to
\begin{mathbox}
\[
\therefore\ c(\theta) = -\log\int \exp\{a(y)b(\theta)+d(y)\}\,dy
\]
\end{mathbox}
Equivalently, moving the negative sign across gives
\begin{mathbox}
\[
-c(\theta) = \log\int \exp\{a(y)b(\theta)+d(y)\}\,dy
\]
\end{mathbox}
and exponentiating both sides recovers the normalisation identity stated above:
\begin{mathbox}
\[
\exp\{-c(\theta)\} = \int \exp\{a(y)b(\theta)+d(y)\}\,dy
\]
\end{mathbox}
This confirms that, given $a(y)$, $b(\theta)$, and $d(y)$, the function $c(\theta)$ is not an independent modelling choice but is fully determined by the requirement that the density integrate to one. $\blacksquare$

\section{Why the Exponential Family Matters}

\subsection*{Toward a Unified Framework}
Nelder and Wedderburn's insight, in 1972, was that a wide range of these non-normal response distributions, Binomial, Poisson, Gamma, and Normal among them, could all be written in the same one-parameter exponential family form introduced above. Once a distribution is expressed that way, the same general machinery (score equations, Fisher information, iteratively reweighted least squares) can be used to estimate its parameters, regardless of which member of the family is in play. The generalised linear model is precisely this: a linear predictor for the mean, exactly as in the classical linear model, combined with a chosen distribution from the exponential family and a link function that connects the mean of that distribution to the linear predictor, allowing the connection to be nonlinear (such as the log-odds link mentioned above) where the data demand it.

This is why the exponential family matters so much to GLM theory. It is not simply a convenient algebraic form, it is the structural bridge that lets the well-understood machinery of the linear model be reused, almost unchanged in spirit, across an entire class of response distributions that would otherwise each require their own separate theory \citep{nelder1972,dobson2018}.

\Reflection
\tbc

\Application
\tbc


\chapter{Building the Exponential Family}

\Exploration
\tbc

\Cognition

\section{Normal Distribution}
The Normal distribution, also known as the Gaussian distribution, is a continuous probability distribution for a real-valued random variable. It occupies a central place in statistics, largely because of the Central Limit Theorem: the average of many statistically independent observations of a random variable, provided that variable has a finite mean and variance, converges to a normal distribution as the number of observations grows large. Consequently, physical quantities that arise as the sum of many independent processes, such as height or measurement error, tend to be approximately normally distributed \citep{casella2002}.

\subsection*{Probability Density Function}
\begin{mathbox}
\[
f(y;\mu) = \frac{1}{(2\pi\sigma^2)^{1/2}}\exp\left\{-\frac{(y-\mu)^2}{2\sigma^2}\right\}
\]
\end{mathbox}

\subsection*{Rewriting in One-Parameter Exponential Form}
Taking logs of both sides gives
\begin{mathbox}
\[
\log f(y;\mu) = -\frac{1}{2}\log(2\pi\sigma^2) - \frac{(y-\mu)^2}{2\sigma^2}
\]
\end{mathbox}
Expanding the square in the exponent separates the terms depending on $y$ alone, $\mu$ alone, and both together:
\begin{mathbox}
\[
\log f(y;\mu) = -\frac{1}{2}\log(2\pi\sigma^2) - \frac{y^2}{2\sigma^2} + \frac{y\mu}{\sigma^2} - \frac{\mu^2}{2\sigma^2}
\]
\end{mathbox}
Exponentiating both sides and grouping terms according to the one-parameter exponential family form $f(y;\theta)=\exp\{a(y)b(\theta)+c(\theta)+d(y)\}$ then gives
\begin{mathbox}
\[
f(y;\mu) = \exp\left\{y\cdot\frac{\mu}{\sigma^2} - \frac{\mu^2}{2\sigma^2} - \left[\frac{1}{2}\log(2\pi\sigma^2)+\frac{y^2}{2\sigma^2}\right]\right\}
\]
\end{mathbox}
so that, with $\theta=\mu$,
\begin{mathbox}
\[
a(y)=y,\quad b(\mu)=\frac{\mu}{\sigma^2},\quad c(\mu)=-\frac{\mu^2}{2\sigma^2},\quad d(y)=-\left[\frac{1}{2}\log(2\pi\sigma^2)+\frac{y^2}{2\sigma^2}\right]
\]
\end{mathbox}
Since $a(y)=y$, the Normal distribution is in canonical form. When the dispersion is treated as fixed, as in the one-parameter form used throughout this portfolio, so that $\sigma^2=1$, $b(\mu)$ reduces to $\mu$ itself, meaning the natural parameter coincides with the mean. The canonical link for the Normal distribution is therefore the identity link $g(\mu)=\mu$, which is consistent with how the mean is modelled directly in the classical linear model discussed in Part III.

\subsection*{Proof: Normalising Term $c(\mu)$}
The value of $c(\mu)$ stated above is not an arbitrary choice; it follows from the general normalisation identity proved above, applied here with $a(y)=y$, $b(\mu)=\mu/\sigma^2$, and $d(y)=-\left[\frac{1}{2}\log(2\pi\sigma^2)+y^2/(2\sigma^2)\right]$.

Starting from the identity $\exp\{-c(\mu)\} = \int\exp\{a(y)b(\mu)+d(y)\}\,dy$ and substituting these terms gives
\begin{mathbox}
\[
\exp\{-c(\mu)\} = \int \exp\left\{\frac{y\mu}{\sigma^2}-\frac{y^2}{2\sigma^2}-\frac{1}{2}\log(2\pi\sigma^2)\right\}dy
\]
\end{mathbox}
Since the last term does not depend on $y$, it can be pulled outside the integral:
\begin{mathbox}
\[
\exp\{-c(\mu)\} = \frac{1}{\sqrt{2\pi\sigma^2}}\int \exp\left\{\frac{y\mu}{\sigma^2}-\frac{y^2}{2\sigma^2}\right\}dy
\]
\end{mathbox}
Completing the square in the remaining exponent, since $\dfrac{y\mu}{\sigma^2}-\dfrac{y^2}{2\sigma^2} = -\dfrac{(y-\mu)^2}{2\sigma^2}+\dfrac{\mu^2}{2\sigma^2}$, gives
\begin{mathbox}
\[
\exp\{-c(\mu)\} = \frac{1}{\sqrt{2\pi\sigma^2}}\exp\left\{\frac{\mu^2}{2\sigma^2}\right\}\int\exp\left\{-\frac{(y-\mu)^2}{2\sigma^2}\right\}dy
\]
\end{mathbox}
The remaining integral is the standard Gaussian integral, which evaluates to $\sqrt{2\pi\sigma^2}$, so this reduces to
\begin{mathbox}
\[
\exp\{-c(\mu)\} = \frac{1}{\sqrt{2\pi\sigma^2}}\cdot\exp\left\{\frac{\mu^2}{2\sigma^2}\right\}\cdot\sqrt{2\pi\sigma^2} = \exp\left\{\frac{\mu^2}{2\sigma^2}\right\}
\]
\end{mathbox}
Equating exponents and solving for $c(\mu)$ therefore gives
\begin{mathbox}
\[
-c(\mu) = \frac{\mu^2}{2\sigma^2} \implies \therefore\ c(\mu) = -\frac{\mu^2}{2\sigma^2}\quad \blacksquare
\]
\end{mathbox}
This confirms that $c(\mu)$ is forced by the same normalisation constraint established generally above, worked through here with the specific $a(y)$ and $d(y)$ that arise from the Normal density.

\section{Binomial Distribution}
The Binomial distribution is a discrete probability distribution for the number of successes in a sequence of $n$ independent trials, each with a boolean outcome, success or failure, where success occurs with probability $p$ and failure with probability $1-p$, since the two outcomes must sum to 1. The special case $n=1$ is the Bernoulli distribution, so the Binomial can be thought of as a generalisation of a sequence of Bernoulli trials \citep{casella2002}.

\subsection*{Probability Mass Function}
\begin{mathbox}
\[
f(y;n,p) = \binom{n}{y}p^y(1-p)^{n-y}, \qquad y=0,1,\dots,n
\]
\end{mathbox}

\subsection*{Rewriting in One-Parameter Exponential Form}
Taking logs of both sides gives
\begin{mathbox}
\[
\log f(y;n,p) = \log\binom{n}{y} + y\log(p) + (n-y)\log(1-p)
\]
\end{mathbox}
Expanding the last term and collecting everything that multiplies $y$ separates the exponent into a piece linear in $y$ and a piece that does not depend on $y$ at all:
\begin{mathbox}
\[
\log f(y;n,p) = y\log(p) - y\log(1-p) + n\log(1-p) + \log\binom{n}{y} = y\log\left(\frac{p}{1-p}\right) + n\log(1-p) + \log\binom{n}{y}
\]
\end{mathbox}
Exponentiating both sides and grouping terms according to $f(y;\theta) = \exp\{a(y)b(\theta)+c(\theta)+d(y)\}$ then gives
\begin{mathbox}
\[
f(y;n,p) = \exp\left\{y\log\left(\frac{p}{1-p}\right) + n\log(1-p) + \log\binom{n}{y}\right\}
\]
\end{mathbox}
so that, with $\theta=p$,
\begin{mathbox}
\[
a(y)=y,\quad b(p)=\log\left(\frac{p}{1-p}\right),\quad c(p) = n\log(1-p),\quad d(y)=\log\binom{n}{y}
\]
\end{mathbox}
Since $a(y)=y$, the Binomial distribution is in canonical form, and $b(p)$ is its canonical parameter. Because $b(p) = \log\left(\frac{p}{1-p}\right)$ is exactly the logit function defined in Part II, the canonical link for the Binomial distribution is the logit link, which is why logistic regression models the log-odds of success as a linear predictor rather than the probability itself.

\subsection*{Proof: Normalising Term $c(p)$}
The value of $c(p)$ stated above is not an arbitrary choice; it follows from the general normalisation identity proved above, applied here with $a(y)=y$, $b(p)=\log\left(\frac{p}{1-p}\right)$, and $d(y)=\log\binom{n}{y}$. Since the Binomial distribution is discrete, the identity holds with a sum in place of an integral.

Starting from $\exp\{-c(p)\} = \sum_{y=0}^n \exp\{a(y)b(p)+d(y)\}$ and substituting these terms gives
\begin{mathbox}
\[
\exp\{-n\log(1-p)\} = \sum_{y=0}^n \exp\left\{y\log\left(\frac{p}{1-p}\right) + \log\binom{n}{y}\right\}
\]
\end{mathbox}
Evaluating the exponent on the right, using $\exp\{A+B\}=\exp\{A\}\exp\{B\}$,
\begin{mathbox}
\[
\exp\left\{y\log\left(\frac{p}{1-p}\right) + \log\binom{n}{y}\right\} = \exp\left\{y\log\left(\frac{p}{1-p}\right)\right\}\cdot\exp\left\{\log\binom{n}{y}\right\} = \left(\frac{p}{1-p}\right)^y\binom{n}{y}
\]
\end{mathbox}
so the right-hand side becomes
\begin{mathbox}
\[
\sum_{y=0}^n \left(\frac{p}{1-p}\right)^y\binom{n}{y} = \left(1+\frac{p}{1-p}\right)^n
\]
\end{mathbox}
by the Binomial theorem. Simplifying the base,
\begin{mathbox}
\[
\left(1+\frac{p}{1-p}\right)^n = \left(\frac{1-p+p}{1-p}\right)^n = \left(\frac{1}{1-p}\right)^n = (1-p)^{-n}
\]
\end{mathbox}
This matches the left-hand side exactly, since
\begin{mathbox}
\[
\exp\{-n\log(1-p)\} = (1-p)^{-n}
\]
\end{mathbox}
Equating both sides and taking logs confirms the identity holds:
\begin{mathbox}
\[
\log\left[\exp\{-n\log(1-p)\}\right] = \log\left[(1-p)^{-n}\right] \implies -n\log(1-p) = -n\log(1-p) \quad\blacksquare
\]
\end{mathbox}
This proves that $c(p) = n\log(1-p)$ is the function that makes $\sum_{y=0}^n f(y;n,p) = 1$.

\section{Poisson Distribution}
The Poisson distribution is a discrete probability distribution that expresses the probability of a number of events occurring in a fixed interval of time (or space), given that these events occur with a known constant mean rate and independently of the time since the last event \citep{casella2002}.

\subsection*{Probability Mass Function}
\begin{mathbox}
\[
f(y;\lambda) = \frac{\lambda^y e^{-\lambda}}{y!}, \qquad \lambda = \text{rate parameter},\ y=0,1,2,\dots
\]
\end{mathbox}

\subsection*{Rewriting in One-Parameter Exponential Form}
Taking logs of both sides gives
\begin{mathbox}
\[
\log f(y;\lambda) = y\log(\lambda) - \lambda - \log(y!)
\]
\end{mathbox}
This is already separated into a piece linear in $y$, a piece depending on $\lambda$ alone, and a piece depending on $y$ alone, so exponentiating and grouping terms according to $f(y;\theta)=\exp\{a(y)b(\theta)+c(\theta)+d(y)\}$ gives
\begin{mathbox}
\[
f(y;\lambda) = \exp\left\{y\log(\lambda) - \lambda - \log(y!)\right\}
\]
\end{mathbox}
so that, with $\theta=\lambda$,
\begin{mathbox}
\[
a(y)=y,\quad b(\lambda)=\log(\lambda),\quad c(\lambda) = -\lambda,\quad d(y) = -\log(y!)
\]
\end{mathbox}
Since $a(y)=y$, the Poisson distribution is in canonical form, and $b(\lambda)=\log(\lambda)$ is its canonical parameter. The canonical link for the Poisson distribution is therefore the log link, $g(\mu)=\log(\mu)$, which is why Poisson regression models the log of the mean count as a linear predictor.

\subsection*{Proof: Normalising Term $c(\lambda)$}
The value of $c(\lambda)$ stated above is not an arbitrary choice; it follows from the general normalisation identity proved above, applied here with $a(y)=y$, $b(\lambda)=\log(\lambda)$, and $d(y)=-\log(y!)$. Since the Poisson distribution is discrete with unbounded support $y=0,1,2,\dots$, the identity holds with an infinite sum in place of an integral.

Starting from $\exp\{-c(\lambda)\} = \sum_{y=0}^{\infty}\exp\{a(y)b(\lambda)+d(y)\}$ and substituting these terms gives
\begin{mathbox}
\[
\exp\{-(-\lambda)\} = \sum_{y=0}^{\infty}\exp\left\{y\log(\lambda) - \log(y!)\right\}
\]
\end{mathbox}
Decomposing the exponent on the right using exponential laws,
\begin{mathbox}
\[
\exp\left\{y\log(\lambda) - \log(y!)\right\} = \frac{\exp\{y\log(\lambda)\}}{\exp\{\log(y!)\}} = \frac{\lambda^y}{y!}
\]
\end{mathbox}
so the identity becomes
\begin{mathbox}
\[
\exp\{\lambda\} = \sum_{y=0}^{\infty}\frac{\lambda^y}{y!}
\]
\end{mathbox}
The right-hand side is exactly the Maclaurin series expansion of the exponential function, which is known to equal $e^\lambda$:
\begin{mathbox}
\[
\sum_{y=0}^{\infty}\frac{\lambda^y}{y!} = e^\lambda = \exp\{\lambda\}
\]
\end{mathbox}
Equating both sides confirms the identity holds,
\begin{mathbox}
\[
\exp\{\lambda\} = \exp\{\lambda\} \quad\blacksquare
\]
\end{mathbox}
This proves that $c(\lambda) = -\lambda$ is the function that makes $\sum_{y=0}^{\infty} f(y;\lambda) = 1$.

\section{Gamma Distribution}
The Gamma distribution is a two-parameter family of continuous probability distributions with many applications across fields such as econometrics and Bayesian statistics \citep{casella2002}. In econometrics, the $(\alpha,\theta)$ shape-scale parameterisation is commonly used to model waiting times, such as time until failure, where it often takes the form of a special case of the Gamma distribution, the Erlang distribution, for integer $\alpha$. In Bayesian statistics, the parameterisation takes the form $(\alpha,\beta)$, with $\beta=1/\theta$ the rate parameter, and the Gamma distribution is widely used as a conjugate prior for several inverse scale parameters, which facilitates analytical tractability in posterior computations. In GLMs, it is more convenient to work with a mean-shape parameterisation, writing the Gamma distribution in terms of its mean $\mu$ and shape $\nu$, with $\nu$ treated as a nuisance parameter \citep{dobson2018}.

\subsection*{Probability Density Function}
\begin{mathbox}
\[
f(y;\mu) = \frac{1}{\Gamma(\nu)}\,y^{-1}\left(\frac{\nu y}{\mu}\right)^\nu\exp\left\{-\frac{\nu y}{\mu}\right\}
\]
\end{mathbox}

\subsection*{Rewriting in One-Parameter Exponential Form}
Taking logs of both sides gives
\begin{mathbox}
\[
\log f(y;\mu) = -\log\Gamma(\nu) - \log(y) + \nu\left[\log(\nu y) - \log(\mu)\right] - \frac{\nu y}{\mu}
\]
\end{mathbox}
Expanding $\log(\nu y)=\log(\nu)+\log(y)$ and collecting terms separates the exponent into a piece linear in $y$, a piece depending on $\mu$ alone, and a piece depending on $y$ alone:
\begin{mathbox}
\[
\log f(y;\mu) = -\frac{\nu y}{\mu} - \nu\log(\mu) + \left[(\nu-1)\log(y) + \nu\log(\nu) - \log\Gamma(\nu)\right]
\]
\end{mathbox}
Exponentiating both sides and grouping terms according to $f(y;\theta) = \exp\{a(y)b(\theta)+c(\theta)+d(y)\}$ then gives
\begin{mathbox}
\[
f(y;\mu) = \exp\left\{y\cdot\left(-\frac{\nu}{\mu}\right) - \nu\log(\mu) + \left[(\nu-1)\log(y) + \nu\log(\nu) - \log\Gamma(\nu)\right]\right\}
\]
\end{mathbox}
so that, with $\theta=\mu$,
\begin{mathbox}
\[
a(y)=y,\quad b(\mu) = -\frac{\nu}{\mu},\quad c(\mu) = -\nu\log(\mu),\quad d(y) = (\nu-1)\log(y)+\nu\log(\nu)-\log\Gamma(\nu)
\]
\end{mathbox}
Since $a(y)=y$, the Gamma distribution is in canonical form, and $b(\mu)=-\nu/\mu$ is its canonical parameter. Because $b(\mu)$ is proportional to $-1/\mu$, the canonical link for the Gamma distribution is the negative reciprocal link, $g(\mu)=-1/\mu$, in contrast to the identity and log links seen for the Normal and Poisson distributions above.

\subsection*{Proof: Normalising Term $c(\mu)$}
The value of $c(\mu)$ stated above is not an arbitrary choice; it follows from the general normalisation identity proved above, applied here with $a(y)=y$, $b(\mu)=-\nu/\mu$, and $d(y)=(\nu-1)\log(y)+\nu\log(\nu)-\log\Gamma(\nu)$.

Starting from $\exp\{-c(\mu)\} = \int_0^\infty \exp\{a(y)b(\mu)+d(y)\}\,dy$ and substituting these terms gives
\begin{mathbox}
\[
\exp\{\nu\log(\mu)\} = \int_0^\infty \exp\left\{-\frac{\nu y}{\mu} + (\nu-1)\log(y) + \nu\log(\nu) - \log\Gamma(\nu)\right\}dy
\]
\end{mathbox}
Decomposing the exponent on the right using exponential laws separates out the part of the integrand that depends on $y$ from the part that does not:
\begin{mathbox}
\[
\exp\left\{-\frac{\nu y}{\mu}\right\} \times \exp\left\{(\nu-1)\log(y)+\nu\log(\nu)-\log\Gamma(\nu)\right\} = \exp\left\{-\frac{\nu y}{\mu}\right\}\times\frac{y^{\nu-1}\nu^\nu}{\Gamma(\nu)}
\]
\end{mathbox}
Plugging this back into the integral and pulling the constants $\nu^\nu/\Gamma(\nu)$, which do not depend on $y$, outside it,
\begin{mathbox}
\[
\int_0^\infty \frac{\nu^\nu}{\Gamma(\nu)}y^{\nu-1}\exp\left\{-\frac{\nu y}{\mu}\right\}dy = \frac{\nu^\nu}{\Gamma(\nu)}\int_0^\infty y^{\nu-1}\exp\left\{-\frac{\nu y}{\mu}\right\}dy
\]
\end{mathbox}
The remaining integral is Euler's integral of the second kind, which states that for $a,b>0$,
\begin{mathbox}
\[
\int_0^\infty y^{a-1}e^{-by}\,dy = \frac{\Gamma(a)}{b^a}
\]
\end{mathbox}
Applying this with $a=\nu$ and $b=\nu/\mu$ gives
\begin{mathbox}
\[
\frac{\nu^\nu}{\Gamma(\nu)}\int_0^\infty y^{\nu-1}\exp\left\{-\frac{\nu y}{\mu}\right\}dy = \frac{\nu^\nu}{\Gamma(\nu)}\cdot\frac{\Gamma(\nu)}{(\nu/\mu)^\nu} = \frac{1}{\mu^{-\nu}} = \mu^\nu
\]
\end{mathbox}
This matches the left-hand side exactly, since $\exp\{\nu\log(\mu)\} = \mu^\nu$. Equating both sides confirms the identity holds:
\begin{mathbox}
\[
\exp\{\nu\log(\mu)\} = \mu^\nu \quad\blacksquare
\]
\end{mathbox}
This proves that $c(\mu) = -\nu\log(\mu)$ is the function that makes $\int_0^\infty f(y;\mu)\,dy = 1$.

\Reflection
\tbc

\Application
\tbc


\chapter{Properties of the Exponential Family}

\Exploration
\tbc

\Cognition

Having built up the Normal, Binomial, Poisson, and Gamma distributions in one-parameter exponential family form, it is worth stepping back and asking what general properties fall out of that form itself, independent of which specific distribution is being used. Three results turn out to matter for everything that follows: the mean and variance of the sufficient statistic $a(Y)$ can be read off directly from $b(\theta)$ and $c(\theta)$ without ever touching an integral, and the score function, the quantity that drives maximum likelihood estimation, has an expected value of exactly zero at the true parameter.

\section{Mean of \texorpdfstring{$a(Y)$}{a(Y)}}
Recall the one-parameter exponential family density from the previous chapter,
\begin{mathbox}
\[
f(y;\theta) = \exp\{a(y)b(\theta)+c(\theta)+d(y)\}.
\]
\end{mathbox}
To find $\mathbb{E}(Y)$ directly would normally require evaluating $\sum_y y f(y)$ or $\int y f(y)\,dy$, which can be difficult or intractable depending on the distribution. The exponential family form lets this be sidestepped entirely by working with $a(Y)$ instead: since $a(Y)$ is a sufficient statistic for $\theta$, it is the natural quantity to study, and when the distribution is in canonical form ($a(y)=y$), a general result for $\mathbb{E}[a(Y)]$ is at the same time a general result for $\mathbb{E}(Y)$.

\textit{Why not just differentiate $c(\theta)$?} It is tempting to think that, since differentiating the normalising term $c(\theta)$ will turn out to produce the mean of $a(Y)$, this could be extracted and differentiated directly. The reason this shortcut does not work on its own is that the derivative of $c(\theta)$ alone is only a function of $\theta$; it says nothing by itself. It is only once it is placed back inside the full density, and the identity that every density must integrate to 1 is used, that a relationship between $c(\theta)$ and the moments of $a(Y)$ emerges. That identity is the actual engine of the derivation below.

\begin{theorem}[Mean of the Sufficient Statistic]
\begin{mathbox}
\[
\mathbb{E}[a(Y)] = -\frac{c'(\theta)}{b'(\theta)}
\]
\end{mathbox}
\end{theorem}
\begin{proof}
Every valid density satisfies
\begin{mathbox}
\[
\int f(y;\theta)\,dy = 1 \quad\text{for all }\theta.
\]
\end{mathbox}
Differentiating both sides with respect to $\theta$, and, under standard regularity conditions, interchanging the order of differentiation and integration,
\begin{mathbox}
\[
\int \frac{\partial f(y;\theta)}{\partial\theta}\,dy = 0.
\]
\end{mathbox}
Differentiating the exponential family density itself with respect to $\theta$, using the chain rule on the exponent,
\begin{mathbox}
\[
\frac{\partial f(y;\theta)}{\partial\theta} = f(y;\theta)\left[a(y)b'(\theta) + c'(\theta)\right].
\]
\end{mathbox}
Substituting this back in and splitting the integral, pulling the $\theta$-only terms $b'(\theta)$ and $c'(\theta)$ outside,
\begin{mathbox}
\[
b'(\theta)\int a(y)f(y;\theta)\,dy + c'(\theta)\int f(y;\theta)\,dy = 0.
\]
\end{mathbox}
The first integral is $\mathbb{E}[a(Y)]$ by definition, and the second is 1, so $b'(\theta)\,\mathbb{E}[a(Y)]+c'(\theta) = 0$, giving
\begin{mathbox}
\[
\mathbb{E}[a(Y)] = -\frac{c'(\theta)}{b'(\theta)}
\]
\end{mathbox}
This is the relationship referred to above: $c'(\theta)$ alone is meaningless without $b'(\theta)$ to divide it by, but together they give the exact mean, with no integration required once $b(\theta)$ and $c(\theta)$ are known.
\end{proof}

\section{Variance of \texorpdfstring{$a(Y)$}{a(Y)}}
To find the variance, the same identity is differentiated a second time with respect to $\theta$:
\begin{mathbox}
\[
\int \frac{\partial^2 f(y;\theta)}{\partial\theta^2}\,dy = 0.
\]
\end{mathbox}
From the mean derivation above, $\dfrac{\partial f}{\partial\theta} = f(y;\theta)\cdot v(\theta,y)$, where $v(\theta,y) = a(y)b'(\theta)+c'(\theta)$. Writing $u=f(y;\theta)$ so that $u'=uv$, the product rule gives $\dfrac{\partial^2 f}{\partial\theta^2} = u'v+uv' = u(v^2+v')$, that is,
\begin{mathbox}
\[
\frac{\partial^2 f}{\partial\theta^2} = f(y;\theta)\left[(a(y)b'(\theta)+c'(\theta))^2 + a(y)b''(\theta)+c''(\theta)\right].
\]
\end{mathbox}
Substituting this into the zero-integral identity above and splitting into two pieces,
\begin{mathbox}
\[
\underbrace{\int f(y;\theta)\left[a(y)b'(\theta)+c'(\theta)\right]^2dy}_{(I)} + \underbrace{\int f(y;\theta)\left[a(y)b''(\theta)+c''(\theta)\right]dy}_{(II)} = 0.
\]
\end{mathbox}

\subsection*{Integral (I)}
Using the mean result above, $c'(\theta) = -b'(\theta)\mathbb{E}[a(Y)]$, so the bracket simplifies to $a(y)b'(\theta)+c'(\theta) = b'(\theta)\left(a(y)-\mathbb{E}[a(Y)]\right)$. Therefore
\begin{mathbox}
\[
(I) = [b'(\theta)]^2\int f(y;\theta)\left(a(y)-\mathbb{E}[a(Y)]\right)^2dy = [b'(\theta)]^2\operatorname{Var}[a(Y)],
\]
\end{mathbox}
since $\mathbb{E}\left[(a(Y)-\mathbb{E}[a(Y)])^2\right]$ is precisely the definition of $\operatorname{Var}[a(Y)]$.

\subsection*{Integral (II)}
Splitting and pulling out the $\theta$-only terms exactly as before,
\begin{mathbox}
\[
(II) = b''(\theta)\,\mathbb{E}[a(Y)] + c''(\theta).
\]
\end{mathbox}

\subsection*{Combining the Two Integrals}
Putting (I) and (II) together,
\begin{mathbox}
\[
[b'(\theta)]^2\operatorname{Var}[a(Y)] + b''(\theta)\,\mathbb{E}[a(Y)] + c''(\theta) = 0.
\]
\end{mathbox}
Substituting the mean result $\mathbb{E}[a(Y)] = -c'(\theta)/b'(\theta)$ and isolating $\operatorname{Var}[a(Y)]$ gives
\begin{mathbox}
\[
\operatorname{Var}[a(Y)] = \frac{b''(\theta)\,c'(\theta) - b'(\theta)\,c''(\theta)}{[b'(\theta)]^3}
\]
\end{mathbox}

\section{Score Function}
The score is the first derivative of the log-likelihood with respect to the parameter of interest,
\begin{mathbox}
\[
U(\theta) = \frac{\partial}{\partial\theta}\ell(\theta) = \frac{\partial}{\partial\theta}\ln L(\theta),
\]
\end{mathbox}
where $L(\theta)$ is the likelihood and $\ell(\theta)=\ln L(\theta)$ is the log-likelihood.
\begin{itemize}
\item If $U(\theta)>0$, the log-likelihood is increasing at $\theta$.
\item If $U(\theta)<0$, the log-likelihood is decreasing at $\theta$.
\item If $U(\theta)=0$, the log-likelihood is flat (stationary) at $\theta$; this is how the maximum likelihood estimate is characterised.
\end{itemize}

\textbf{An important property: $\mathbb{E}[U(\theta)]=0$.} A zero expected score reflects the fact that the true parameter sits at the centre of the information provided by the data: on average, the log-likelihood has no tendency to push away from the true $\theta$. This holds for any regular density, not only the exponential family: since
\begin{mathbox}
\[
U(\theta) = \frac{1}{f(y;\theta)}\frac{\partial f(y;\theta)}{\partial\theta},
\]
\end{mathbox}
\begin{mathbox}
\[
\mathbb{E}[U(\theta)] = \int f(y;\theta)\,U(\theta)\,dy = \int \frac{\partial f(y;\theta)}{\partial\theta}\,dy = 0,
\]
\end{mathbox}
directly from the zero-integral identity used throughout this chapter. The exponential family proof below is the same identity, made concrete using the explicit form of $f(y;\theta)$.

\begin{theorem}[Zero Expected Score]
$\mathbb{E}[U(\theta)] = 0$ for any regular density, exponential family or otherwise.
\end{theorem}
\begin{proof}
For $f(y;\theta)=\exp[a(y)b(\theta)+c(\theta)+d(y)]$, the log-density is linear in $b(\theta)$ and $c(\theta)$, so the score is
\begin{mathbox}
\[
U(\theta) = \frac{\partial}{\partial\theta}\log f(Y;\theta) = a(Y)b'(\theta)+c'(\theta).
\]
\end{mathbox}
Taking the expectation, using linearity, and substituting the mean result $\mathbb{E}[a(Y)]=-c'(\theta)/b'(\theta)$,
\begin{mathbox}
\[
\mathbb{E}[U(\theta)] = b'(\theta)\,\mathbb{E}[a(Y)]+c'(\theta) = b'(\theta)\left(-\frac{c'(\theta)}{b'(\theta)}\right)+c'(\theta) = -c'(\theta)+c'(\theta) = 0.\ \blacksquare
\]
\end{mathbox}
\end{proof}

\section{Fisher's Information}
Fisher Information answers a single question: how much does one observation tell us about the parameter $\theta$? It admits two equivalent definitions, both of which follow directly from the exponential family results already derived above.

\subsection*{As the Variance of the Score}
Fisher Information is defined as the variance of the score,
\begin{mathbox}
\[
\mathcal{I}(\theta) = \operatorname{Var}[U(\theta)].
\]
\end{mathbox}
Because $U(\theta) = a(Y)b'(\theta)+c'(\theta)$ and $c'(\theta)$ does not depend on $Y$, $\operatorname{Var}[U(\theta)] = [b'(\theta)]^2\operatorname{Var}[a(Y)]$. Substituting the variance result derived above,
\begin{mathbox}
\[
\mathcal{I}(\theta) = [b'(\theta)]^2\cdot\frac{b''(\theta)\,c'(\theta)-b'(\theta)\,c''(\theta)}{[b'(\theta)]^3} = \frac{b''(\theta)\,c'(\theta)-b'(\theta)\,c''(\theta)}{b'(\theta)}.
\]
\end{mathbox}

\subsection*{As the Expected Curvature of the Log-Likelihood}
Differentiating the score once more with respect to $\theta$,
\begin{mathbox}
\[
\frac{\partial^2\ell}{\partial\theta^2} = a(Y)b''(\theta)+c''(\theta).
\]
\end{mathbox}
Taking the expectation and substituting $\mathbb{E}[a(Y)]=-c'(\theta)/b'(\theta)$ from the mean result above,
\begin{mathbox}
\[
\mathbb{E}\left[\frac{\partial^2\ell}{\partial\theta^2}\right] = b''(\theta)\,\mathbb{E}[a(Y)]+c''(\theta) = -\frac{b''(\theta)\,c'(\theta)}{b'(\theta)}+c''(\theta) = -\mathcal{I}(\theta).
\]
\end{mathbox}
So the two definitions coincide:
\begin{mathbox}
\[
\mathcal{I}(\theta) = \operatorname{Var}[U(\theta)] = -\mathbb{E}\left[\frac{\partial^2\ell}{\partial\theta^2}\right]
\]
\end{mathbox}
A sharply peaked log-likelihood (high curvature, a large negative second derivative) corresponds to high information; a flat log-likelihood leaves the parameter poorly determined. For $n$ independent observations, information is additive, $\mathcal{I}_n(\theta) = n\cdot\mathcal{I}_1(\theta)$, and the Cram\'er--Rao bound states that no unbiased estimator can beat
\begin{mathbox}
\[
\operatorname{Var}(\hat\theta) \geq \frac{1}{\mathcal{I}_n(\theta)},
\]
\end{mathbox}
with the maximum likelihood estimator achieving this bound asymptotically.

\Reflection
\tbc

\Application

\textbf{Case study: the Normal distribution.} Rather than differentiating $\log p(x;\mu)$ directly, the Normal case is worked out here from its one-parameter exponential family form, to stay consistent with how every other distribution in this portfolio is handled. From the previous chapter, with $\theta=\mu$ and $\sigma^2$ treated as known,
\begin{mathbox}
\[
a(y)=y,\quad b(\mu)=\frac{\mu}{\sigma^2},\quad c(\mu)=-\frac{\mu^2}{2\sigma^2},
\]
\end{mathbox}
so that
\begin{mathbox}
\[
b'(\mu)=\frac{1}{\sigma^2},\quad b''(\mu)=0,\quad c'(\mu)=-\frac{\mu}{\sigma^2},\quad c''(\mu)=-\frac{1}{\sigma^2}.
\]
\end{mathbox}
Substituting these into the general information formula derived above,
\begin{mathbox}
\[
\mathcal{I}(\mu) = \frac{b''(\mu)\,c'(\mu)-b'(\mu)\,c''(\mu)}{b'(\mu)} = \frac{0-\left(\frac{1}{\sigma^2}\right)\left(-\frac{1}{\sigma^2}\right)}{1/\sigma^2} = \frac{1/\sigma^4}{1/\sigma^2} = \frac{1}{\sigma^2}.
\]
\end{mathbox}
For $n$ independent observations, $\mathcal{I}_n(\mu)=n/\sigma^2$, and the Cram\'er--Rao bound gives $\operatorname{Var}(\hat\mu)\geq \sigma^2/n$.

A useful feature of this case falls straight out of the exponential family form: because $b''(\mu)=0$, the curvature $a(Y)b''(\mu)+c''(\mu)$ reduces to $c''(\mu)$ alone, with no $Y$ in it at all. The observed information therefore equals the expected information exactly, for every sample, not just on average, a structural property of the Normal distribution that does not hold for a distribution such as the Poisson.

A simulation with $\mu=5$ under two variances, $\sigma^2=25$ and $\sigma^2=1$, illustrates these results. Panel 1 shows the log-likelihood sharpening as variance decreases, holding $n$ fixed. Panel 2 shows the score fluctuating more widely across repeated single-observation draws when $\sigma^2$ is small, exactly the larger variance that defines higher Fisher Information. Panel 3 shows the information as flat in $\mu$ for both scenarios, confirming $\mathcal{I}(\mu)$ does not depend on where the true mean sits. Panel 4 tracks the Cram\'er--Rao bound: empirical variance of 4000 simulated MLEs against the theoretical bound $\sigma^2/n$, on a log-log scale, for both variances. Panel 5 plots $\mathcal{I}(\mu)=1/\sigma^2$ against $\sigma^2$ directly, marking the two scenarios and making the inverse relationship visible.

\begin{center}
\textit{[Figure: normal\_fisher\_panels\_5.png -- not embedded in source]}
\end{center}

Numerically, the theoretical information values $\mathcal{I}(\mu)=0.04$ and $\mathcal{I}(\mu)=1$ match a Monte Carlo estimate of $\operatorname{Var}(\hat\mu)$ at $n=100$ closely (0.2435 against a bound of 0.25 for $\sigma^2=25$; 0.0102 against a bound of 0.01 for $\sigma^2=1$): a 25-fold difference in variance produces exactly a 25-fold difference in information.
